% =============================================================================
% ANEXOS — MATERIAL DE REFERENCIA
% =============================================================================
\chapter[Anexos y Material de Referencia]{Anexos y Material de Referencia}%
\label{ch:anexos}
\resumenCapitulo{Este capítulo reúne el material de referencia de consulta rápida que
complementa los procedimientos del manual: una hoja de comandos por fase, el catálogo
completo de \textit{endpoints} REST, las plantillas íntegras de los archivos de
configuración, la matriz extendida de variables, una sección de preguntas frecuentes y las
listas de verificación imprimibles por fase de instalación.}

Los anexos no introducen pasos nuevos: condensan, en formato de referencia, la información
dispersa en los capítulos anteriores para que el operador la tenga a mano durante el
despliegue. Cada anexo remite, mediante referencias cruzadas, al capítulo donde el
procedimiento correspondiente se explica en detalle.\par

% -----------------------------------------------------------------------------
\section{Anexo A — Hoja de Referencia Rápida de Comandos}%
\label{sec:anexo-comandos}

La tabla~\ref{tab:cheatsheet} consolida los comandos esenciales de cada fase del ciclo de
instalación, en el orden en que se ejecutan.\par

\vspace{0.2cm}
\noindent
\renewcommand{\arraystretch}{1.3}
\fontsize{8.3}{10}\selectfont\sffamily
\begin{tabularx}{\dimexpr\textwidth-2pt\relax}{|>{\bfseries\color{colorPrimario}}l|>{\ttfamily\raggedright\arraybackslash}X|}
\hline
\rowcolor{headerTabla}
\sffamily\textbf{Fase} & \sffamily\textbf{Comando} \\
\hline
Verificar JDK & java -version \\
\hline
\rowcolor{grisTecnico}
Verificar Node & node -v \&\& npm -v \\
\hline
Verificar Docker & docker compose version \\
\hline
\rowcolor{grisTecnico}
Clonar código & git clone <url>/ethoshubB.git \&\& git clone <url>/ethoshubF.git \\
\hline
Configurar backend & cp .env.example .env \\
\hline
\rowcolor{grisTecnico}
Migrar BD & psql "\$PGURL" -f src/main/resources/db/migration/V1\_\_*.sql \\
\hline
Construir (completo) & ./build.sh -{}-prod \\
\hline
\rowcolor{grisTecnico}
Construir (solo BE) & ./build.sh -{}-skip-fe \\
\hline
Ejecutar JAR & java -jar target/ethoshub-0.0.1-SNAPSHOT.jar -{}-spring.profiles.active=prod \\
\hline
\rowcolor{grisTecnico}
Construir imagen & docker build -t ethoshub:2.0.0 . \\
\hline
Levantar Compose & docker compose up -d -{}-build \\
\hline
\rowcolor{grisTecnico}
Ver logs & docker compose logs -f ethoshub-backend \\
\hline
Salud & curl -s http://localhost:8080/actuator/health \\
\hline
\rowcolor{grisTecnico}
Detener & docker compose stop \\
\hline
Purga total & docker compose down -v \\
\hline
\end{tabularx}\normalfont\normalsize%
\label{tab:cheatsheet}

% -----------------------------------------------------------------------------
\section{Anexo B — Catálogo Completo de Endpoints REST}%
\label{sec:anexo-endpoints}

La tabla~\ref{tab:endpoints-full} cataloga los controladores REST expuestos por el backend,
agrupados por dominio funcional. El operador puede inspeccionarlos en Swagger UI
(apartado~\ref{subsec:swagger}) para confirmar que la API se desplegó por completo.\par

\vspace{0.2cm}
\renewcommand{\arraystretch}{1.3}
\fontsize{8.2}{9.8}\selectfont\sffamily
\setlength{\tabcolsep}{4pt}
\begin{longtable}{|>{\bfseries\color{colorPrimario}\raggedright\arraybackslash}p{2.5cm}|>{\ttfamily\raggedright\arraybackslash}p{5.2cm}|>{\raggedright\arraybackslash}p{7.2cm}|}
\caption{Catálogo de controladores REST del backend, por dominio funcional.}\label{tab:endpoints-full}\\
\hline
\rowcolor{headerTabla}
\sffamily\textbf{Dominio} & \sffamily\textbf{Controlador} & \sffamily\textbf{Responsabilidad} \\
\hline
\endfirsthead
\hline
\rowcolor{headerTabla}
\sffamily\textbf{Dominio} & \sffamily\textbf{Controlador} & \sffamily\textbf{Responsabilidad} \\
\hline
\endhead
\hline
\multicolumn{3}{r}{\sffamily\tiny\color{grisISO}\textit{continúa en la página siguiente}}\\
\endfoot
\hline
\endlastfoot
Auth & AuthController & Registro, login y sincronización OAuth. \\
\hline
\rowcolor{grisTecnico}
Auth & ForgotPasswordController & Recuperación de contraseña por correo. \\
\hline
Seguridad & SecurityController & Cambio de contraseña/email y baja de cuenta con OTP. \\
\hline
\rowcolor{grisTecnico}
Perfil & ProfileController & Perfil profesional básico del usuario. \\
\hline
Proyectos & ProjectController & CRUD de proyectos, subida multi-PDF, Google Drive. \\
\hline
\rowcolor{grisTecnico}
Experiencia & WorkExperienceController & Experiencia laboral con campos geográficos. \\
\hline
Formación & AcademicRecordController & Registros académicos y formación. \\
\hline
\rowcolor{grisTecnico}
Skills & SkillController & Hard y soft skills (con \textit{soft-delete}). \\
\hline
Portafolio & PortfolioSettingsController & Ajustes de visibilidad, SEO y contraseña. \\
\hline
\rowcolor{grisTecnico}
Portafolio & PortfolioPublicController & Vista pública del portafolio. \\
\hline
CV Studio & CvDocumentController & Edición de CV, IA y compilación a PDF. \\
\hline
\rowcolor{grisTecnico}
CV Studio & FileUploadController & Subida de archivos del CV. \\
\hline
Conexiones & ConnectionController & Conexiones a apps externas (URL-first). \\
\hline
\rowcolor{grisTecnico}
Chat & ChatController & Mensajería en tiempo real (Supabase Realtime). \\
\hline
Notificaciones & NotificationController & Notificaciones del usuario. \\
\hline
\rowcolor{grisTecnico}
Preferencias & PreferenceController & Preferencias e idioma del usuario. \\
\hline
Reclutador & TalentController & Descubrimiento de talento y \textit{pipeline} Kanban. \\
\hline
\rowcolor{grisTecnico}
Reclutador & RecruiterProfileController & Perfil del reclutador. \\
\hline
Reclutador & RecruiterAiController & \textit{Insights} de IA y búsqueda NLP. \\
\hline
\rowcolor{grisTecnico}
Admin & AdminMetricsController & Métricas del panel de administración. \\
\hline
Admin & AdminModerationController & Moderación de contenido. \\
\hline
\rowcolor{grisTecnico}
Admin & AdminProfileController & Gestión de perfiles. \\
\hline
Admin & AdminSkillController & Gestión del catálogo de skills. \\
\hline
\rowcolor{grisTecnico}
Admin & AdminEmailController & Envío de correos masivos. \\
\hline
Infra & HealthController & Sonda de salud del servicio. \\
\hline
\end{longtable}

% -----------------------------------------------------------------------------
\section{Anexo C — Plantillas Completas de Archivos de Configuración}%
\label{sec:anexo-plantillas}

Este anexo reproduce, íntegras, las plantillas de los archivos de configuración que el
operador debe rellenar. Sustituya los valores entre \comando{<...>} por los reales.\par

\subsection{C.1 — \texttt{ethoshubB/.env} (backend)}%
\label{subsec:plantilla-env-be}

\begin{lstlisting}[language=bash]
# === Servidor ===============================================
SERVER_PORT=8080
SPRING_PROFILES_ACTIVE=prod

# === Supabase ===============================================
SUPABASE_URL=https://<project-ref>.supabase.co
SUPABASE_ANON_KEY=<anon-key>
SUPABASE_SERVICE_ROLE_KEY=<service-role-key>   # SECRETO

# === Validacion de JWT (defina UNO) =========================
# HS256 (secreto compartido):
SUPABASE_JWT_SECRET=
# RS256 (JWKS, recomendado):
SPRING_SECURITY_OAUTH2_RESOURCESERVER_JWT_JWK_SET_URI=https://<ref>.supabase.co/auth/v1/.well-known/jwks.json

# === IA (Google Gemini) =====================================
GEMINI_API_KEY=<gemini-key>
\end{lstlisting}

\subsection{C.2 — \texttt{ethoshubF/.env.production} (frontend)}%
\label{subsec:plantilla-env-fe}

\begin{lstlisting}[language=bash]
# Same-origin en el monolito: VITE_API_URL vacia -> baseURL '/api'
VITE_API_URL=
VITE_SUPABASE_URL=https://<project-ref>.supabase.co
VITE_SUPABASE_ANON_KEY=<anon-key>

# Google Drive Picker (credenciales publicas de cliente)
VITE_GOOGLE_API_KEY=<google-api-key>
VITE_GOOGLE_CLIENT_ID=<google-client-id>.apps.googleusercontent.com
\end{lstlisting}

\subsection{C.3 — \texttt{docker-compose.yml}}%
\label{subsec:plantilla-compose}

\begin{lstlisting}[language=bash]
services:
  ethoshub-backend:
    build:
      context: .
    ports:
      - "8080:8080"
    environment:
      - SPRING_PROFILES_ACTIVE=prod
      - SUPABASE_DB_URL=jdbc:postgresql://<host>:5432/postgres
      - SUPABASE_DB_USER=postgres
      - SUPABASE_DB_PASSWORD=<password>
      - SUPABASE_JWKS_URI=https://<ref>.supabase.co/auth/v1/.well-known/jwks.json
    volumes:
      - .:/app
      - spring-cache:/root/.m2
    networks:
      - ethoshub-net
volumes:
  spring-cache:
networks:
  ethoshub-net:
    driver: bridge
\end{lstlisting}

\begin{NotaConfiguracion}
Para no exponer secretos en \comando{docker-compose.yml}, sustituya el bloque
\comando{environment} por la directiva \comando{env\_file: .env} y mantenga el archivo
\comando{.env} fuera del control de versiones, según el principio de custodia del
apartado~\ref{subsec:servicerole}.
\end{NotaConfiguracion}

% -----------------------------------------------------------------------------
\section{Anexo D — Matriz Extendida de Variables de Entorno}%
\label{sec:anexo-vars}

La tabla~\ref{tab:vars-ext} amplía la matriz del apartado~\ref{subsec:matriz-variables}
con el valor por defecto y el componente de cada variable.\par

\vspace{0.2cm}
\noindent
\renewcommand{\arraystretch}{1.3}
\fontsize{8.2}{9.8}\selectfont\sffamily
\begin{tabularx}{\dimexpr\textwidth-2pt\relax}{|>{\ttfamily\bfseries\color{colorPrimario}\raggedright\arraybackslash}p{4.6cm}|c|>{\raggedright\arraybackslash}X|}
\hline
\rowcolor{headerTabla}
\sffamily\textbf{Variable} & \sffamily\textbf{Comp.} & \sffamily\textbf{Valor por defecto / nota} \\
\hline
SERVER\_PORT & BE & 8080 \\
\hline
\rowcolor{grisTecnico}
SPRING\_PROFILES\_ACTIVE & BE & local (usar prod en despliegue) \\
\hline
SUPABASE\_URL & BE & — (obligatoria) \\
\hline
\rowcolor{grisTecnico}
SUPABASE\_ANON\_KEY & BE & — (obligatoria) \\
\hline
SUPABASE\_SERVICE\_ROLE\_KEY & BE & — (secreta, obligatoria) \\
\hline
\rowcolor{grisTecnico}
...JWT\_JWK\_SET\_URI & BE & JWKS de Supabase \\
\hline
GEMINI\_API\_KEY & BE & — (IA; sin ella, IA devuelve error) \\
\hline
\rowcolor{grisTecnico}
APP\_BASE\_URL & BE & http://localhost:8080 \\
\hline
APP\_FRONTEND\_URL & BE & http://localhost:5173 \\
\hline
\rowcolor{grisTecnico}
VITE\_API\_URL & FE & vacía (same-origin /api) \\
\hline
VITE\_SUPABASE\_URL & FE & — (obligatoria) \\
\hline
\rowcolor{grisTecnico}
VITE\_SUPABASE\_ANON\_KEY & FE & — (obligatoria) \\
\hline
VITE\_GOOGLE\_CLIENT\_ID & FE & — (Drive Picker, opcional) \\
\hline
\rowcolor{grisTecnico}
VITE\_GOOGLE\_API\_KEY & FE & — (Drive Picker, opcional) \\
\hline
\end{tabularx}\normalfont\normalsize%
\label{tab:vars-ext}

% -----------------------------------------------------------------------------
\section{Anexo E — Preguntas Frecuentes de Instalación (FAQ)}%
\label{sec:anexo-faq}

\textbf{¿Debo instalar PostgreSQL en el servidor?}\par
No. En producción la base de datos la provee Supabase. PostgreSQL solo aparece localmente a
través de Testcontainers durante las pruebas (apartado~\ref{sec:testcontainers}).\par

\textbf{¿Puedo desplegar sin Docker?}\par
Sí. Ejecute directamente \comando{java -jar} \ruta{target/ethoshub-0.0.1-SNAPSHOT.jar}
\ruta{-{}-spring.profiles.active=prod}. Docker es la vía recomendada por su reproducibilidad, no
una exigencia (apartado~\ref{sec:prod}).\par

\textbf{¿Por qué la cámara no funciona en la red local?}\par
Porque el origen es HTTP inseguro. Aplique las excepciones de navegador del
apartado~\ref{sec:pwa-http} o sirva la aplicación bajo HTTPS.\par

\textbf{¿El sistema arranca sin la clave de Gemini?}\par
Sí. Solo se deshabilitan los \textit{endpoints} de IA; el resto opera con normalidad
(apartado~\ref{sec:degradacion-ia}).\par

\textbf{¿Dónde van los archivos compilados?}\par
El \textit{frontend} compila a \comando{dist/}, que se copia a \comando{static/} dentro del
JAR. Los artefactos de LaTeX de \emph{este} manual se generan en \comando{build/}.\par

\textbf{¿Tengo que definir \texttt{VITE\_API\_URL}?}\par
No en el monolito ni en desarrollo (el proxy de Vite la cubre). Se deja vacía para que
Axios use la ruta relativa \comando{/api}.\par

\textbf{¿Qué hago si el puerto 8080 está ocupado?}\par
Identifique el proceso con \comando{ss -ltnp | grep 8080} y deténgalo, o cambie
\comando{SERVER\_PORT} (capítulo~\ref{ch:troubleshooting}).\par

% -----------------------------------------------------------------------------
\section{Anexo F — Listas de Verificación Imprimibles por Fase}%
\label{sec:anexo-checklists}

\textbf{F.1 — Pre-requisitos (capítulo~\ref{ch:requisitos})}\par
\begin{checklist}
    \item JDK 17+ instalado y en el \texttt{PATH}.
    \item Node.js 20+ y npm 10+ disponibles.
    \item Docker Engine 24+ y Compose v2 activos.
    \item Credenciales de Supabase obtenidas (URL, anon, service-role, JWKS).
    \item Clave de Gemini obtenida (si se requieren funciones de IA).
    \item Puerto 8080 libre.
\end{checklist}

\textbf{F.2 — Preparación (capítulo~\ref{ch:preparacion})}\par
\begin{checklist}
    \item \texttt{ethoshubB} y \texttt{ethoshubF} clonados como hermanos.
    \item \texttt{.env} del backend rellenado.
    \item \texttt{.env.production} del frontend rellenado.
    \item Migraciones aplicadas en orden sobre Supabase.
\end{checklist}

\textbf{F.3 — Construcción y despliegue (capítulos~\ref{ch:construccion} y~\ref{ch:despliegue})}\par
\begin{checklist}
    \item \texttt{build.sh} finalizó con éxito y generó el JAR.
    \item \texttt{static/} contiene \texttt{index.html} y \texttt{assets/}.
    \item Imagen Docker construida (si aplica).
    \item Contenedor en estado \texttt{Up} o JAR en ejecución.
\end{checklist}

\textbf{F.4 — Verificación y cierre (capítulos~\ref{ch:verificacion} y~\ref{ch:desinstalacion})}\par
\begin{checklist}
    \item \texttt{/actuator/health} responde \texttt{UP}.
    \item Swagger UI carga los controladores.
    \item 401/403 se devuelven correctamente.
    \item Excepciones de navegador aplicadas (si HTTP) y previstas para reversión.
    \item Procedimiento de purga y limpieza de cliente conocido.
\end{checklist}

% -----------------------------------------------------------------------------
\section{Anexo G — Endurecimiento del Despliegue de Producción}%
\label{sec:anexo-hardening}

El despliegue básico de los capítulos~\ref{ch:despliegue} y~\ref{ch:seguridad} sirve la
aplicación en HTTP por el puerto 8080. Para un entorno de producción real conviene
\textbf{endurecerlo} con tres medidas estándar: terminación TLS mediante proxy inverso,
ejecución como servicio del sistema y restricción del cortafuegos. Este anexo es opcional
pero recomendado.\par

\subsection{G.1 — Proxy Inverso con TLS (Nginx + Certbot)}%
\label{subsec:nginx}

Un proxy inverso Nginx termina TLS y reenvía el tráfico cifrado al monolito en 8080. Esto
sirve la PWA bajo \comando{https://}, lo que \textbf{elimina la necesidad} de las
excepciones de navegador del apartado~\ref{sec:pwa-http}.\par

\begin{lstlisting}[language=bash]
# /etc/nginx/sites-available/ethoshub
server {
    listen 80;
    server_name bytebusters.tis.cs.umss.edu.bo;
    location / {
        proxy_pass http://127.0.0.1:8080;
        proxy_set_header Host              $host;
        proxy_set_header X-Real-IP         $remote_addr;
        proxy_set_header X-Forwarded-For   $proxy_add_x_forwarded_for;
        proxy_set_header X-Forwarded-Proto $scheme;
    }
}
\end{lstlisting}

\begin{lstlisting}[language=bash]
# Habilitar el sitio y emitir el certificado TLS gratuito (Let's Encrypt)
$ sudo ln -s /etc/nginx/sites-available/ethoshub /etc/nginx/sites-enabled/
$ sudo nginx -t && sudo systemctl reload nginx
$ sudo certbot --nginx -d bytebusters.tis.cs.umss.edu.bo
\end{lstlisting}

\begin{VerificacionOK}
Tras emitir el certificado, la aplicación responde en \comando{https://} con un candado
válido. La cámara y el micrófono de la PWA funcionan \textbf{sin} banderas de navegador,
pues el origen ya es seguro. Esta es la solución definitiva recomendada frente a las
excepciones temporales del apartado~\ref{sec:pwa-http}.
\end{VerificacionOK}

\subsection{G.2 — Ejecución como Servicio del Sistema (systemd)}%
\label{subsec:systemd}

Para que el JAR arranque al iniciar el servidor y se reinicie ante fallos, defínalo como
una unidad de \comando{systemd}:\par

\begin{lstlisting}[language=bash]
# /etc/systemd/system/ethoshub.service
[Unit]
Description=EthosHub Monolito
After=network.target

[Service]
User=ethoshub
WorkingDirectory=/opt/ethoshub
EnvironmentFile=/opt/ethoshub/.env
ExecStart=/usr/bin/java -jar /opt/ethoshub/ethoshub.jar --spring.profiles.active=prod
SuccessExitStatus=143
Restart=on-failure
RestartSec=5

[Install]
WantedBy=multi-user.target
\end{lstlisting}

\begin{lstlisting}[language=bash]
$ sudo systemctl daemon-reload
$ sudo systemctl enable --now ethoshub
$ sudo systemctl status ethoshub      # debe figurar "active (running)"
$ journalctl -u ethoshub -f           # logs en vivo
\end{lstlisting}

\subsection{G.3 — Restricción del Cortafuegos (ufw)}%
\label{subsec:firewall}

Exponga únicamente los puertos necesarios. Con el proxy Nginx, el puerto 8080 \textbf{no}
debe ser accesible desde el exterior; solo 80/443.\par

\begin{lstlisting}[language=bash]
$ sudo ufw allow 22/tcp        # SSH
$ sudo ufw allow 80,443/tcp    # HTTP/HTTPS (Nginx)
$ sudo ufw deny 8080/tcp       # bloquea acceso externo directo al JAR
$ sudo ufw enable
$ sudo ufw status numbered
\end{lstlisting}

\begin{AlertaSeguridad}
Verifique que el puerto 8080 \textbf{no} responde desde una máquina externa
(\comando{curl http://<ip-publica>:8080} debe agotar el tiempo). El único punto de entrada
público debe ser el proxy TLS. Dejar 8080 abierto anula la terminación TLS y reexpone la
PWA en HTTP inseguro.
\end{AlertaSeguridad}

% -----------------------------------------------------------------------------
\section{Anexo H — Topología de Despliegue de Producción y Glosario Extendido}%
\label{sec:anexo-topologia}

La figura~\ref{fig:topologia-prod} integra los componentes de un despliegue endurecido: el
navegador accede por HTTPS al proxy Nginx, que reenvía al monolito; este se conecta a los
servicios externos.\par

\vspace{0.3cm}
\begin{figure}[h!]
\centering
\begin{tikzpicture}[node distance=0.8cm and 1.5cm, every node/.style={font=\sffamily\scriptsize}]
    \node[c4persona, minimum width=2.0cm, minimum height=0.9cm] (nav) {Navegador\\\scriptsize PWA};
    \node[c4contenedor, right=of nav, minimum width=2.0cm, minimum height=0.9cm] (nginx) {Nginx\\\scriptsize TLS :443};
    \node[c4sistema, right=of nginx, minimum width=2.2cm, minimum height=0.9cm] (jar) {Monolito\\\scriptsize :8080};
    \node[c4datos, above right=0.6cm and 1.2cm of jar] (sb) {Supabase};
    \node[c4externo, right=1.2cm of jar] (gem) {Gemini};
    \node[c4externo, below right=0.6cm and 1.2cm of jar] (drv) {Drive};
    \draw[c4flecha] (nav) -- node[c4etiqueta, above]{HTTPS} (nginx);
    \draw[c4flecha] (nginx) -- node[c4etiqueta, above]{:8080} (jar);
    \draw[c4flecha] (jar) -- (sb);
    \draw[c4flecha] (jar) -- (gem);
    \draw[c4flecha] (jar) -- (drv);
\end{tikzpicture}
\caption{Topología de un despliegue de producción endurecido: TLS en el borde (Nginx) y monolito interno conectado a los servicios gestionados.}%
\label{fig:topologia-prod}
\end{figure}

\par\vspace{0.2cm}
La tabla~\ref{tab:glosario-ext} amplía el glosario del apartado~\ref{sec:glosario} con
términos de infraestructura empleados en los anexos.\par

\vspace{0.2cm}
\noindent
\renewcommand{\arraystretch}{1.35}
\fontsize{8.7}{10.4}\selectfont\sffamily
\begin{tabularx}{\dimexpr\textwidth-2pt\relax}{|>{\bfseries\color{colorPrimario}}l|X|}
\hline
\rowcolor{headerTabla}
\textbf{Término} & \textbf{Definición} \\
\hline
Proxy inverso & Servidor (p.\,ej. Nginx) que recibe peticiones del exterior y las reenvía a un servicio interno, terminando TLS. \\
\hline
\rowcolor{grisTecnico}
Terminación TLS & Punto donde el tráfico cifrado HTTPS se descifra antes de reenviarse en claro a la red interna. \\
\hline
systemd & Sistema de inicio de Linux que gestiona servicios; permite arranque automático y reinicio ante fallos. \\
\hline
\rowcolor{grisTecnico}
\textit{Pooler} & Agrupador de conexiones de Supabase que multiplexa conexiones a PostgreSQL (puerto 5432). \\
\hline
Same-origin & Política por la que cliente y API comparten esquema, host y puerto; el monolito la cumple por diseño. \\
\hline
\rowcolor{grisTecnico}
Smoke test & Prueba mínima que confirma que el servicio responde, sin validar toda su funcionalidad. \\
\hline
\end{tabularx}\normalfont\normalsize%
\label{tab:glosario-ext}

% -----------------------------------------------------------------------------
\section{Anexo I — Referencia Anotada de \texttt{application.yml}}%
\label{sec:anexo-appyml}

El backend lee su configuración de \comando{application.yml} (perfil base) y
\comando{application-prod.yml} (producción). El listado siguiente anota las claves más
relevantes y su variable de entorno asociada. Los valores entre \comando{\$\{...\}} se
resuelven desde el entorno; el texto tras los dos puntos es el valor por defecto.\par

\begin{lstlisting}[language=bash]
server:
  port: ${SERVER_PORT:8080}            # Puerto HTTP del monolito

spring:
  profiles:
    active: ${SPRING_PROFILES_ACTIVE:local}
  datasource:
    url: ${DATABASE_URL}               # JDBC del pooler de Supabase
    username: ${DATABASE_USERNAME}
    password: ${DATABASE_PASSWORD}
  security:
    oauth2:
      resourceserver:
        jwt:
          jwk-set-uri: ${SUPABASE_JWK_SET_URI}   # Valida firmas JWT
  jackson:
    default-property-inclusion: non_null         # Omite nulos en las respuestas
    time-zone: UTC                                # Fechas en UTC

supabase:
  url: ${SUPABASE_URL}
  anon-key: ${SUPABASE_ANON_KEY}
  service-role-key: ${SUPABASE_SERVICE_ROLE_KEY}  # SECRETO

gemini:
  api-key: ${GEMINI_API_KEY}            # IA: CV Studio y Recruiter Insights

server.tomcat:
  max-http-form-post-size: -1           # Sin limite (correo masivo de admin)
  max-swallow-size: -1
\end{lstlisting}

\begin{NotaConfiguracion}
La configuración \comando{max-http-form-post-size: -1} y \comando{max-swallow-size: -1} de
Tomcat elimina el límite del tamaño de formulario, necesario para el envío de correo masivo
desde el panel de administración. No la modifique salvo que su política de seguridad exija
acotar el tamaño de las peticiones.
\end{NotaConfiguracion}

% -----------------------------------------------------------------------------
\section{Anexo J — Ajuste de Rendimiento de la JVM y Observabilidad}%
\label{sec:anexo-jvm}

Para entornos de producción con memoria acotada, conviene fijar los parámetros de la JVM en
lugar de depender de los valores automáticos. La tabla~\ref{tab:jvm} resume las opciones
habituales.\par

\vspace{0.2cm}
\noindent
\renewcommand{\arraystretch}{1.3}
\fontsize{8.5}{10.2}\selectfont\sffamily
\begin{tabularx}{\dimexpr\textwidth-2pt\relax}{|>{\ttfamily\bfseries\color{colorPrimario}\raggedright\arraybackslash}p{5.0cm}|X|}
\hline
\rowcolor{headerTabla}
\sffamily\textbf{Opción JVM} & \sffamily\textbf{Efecto} \\
\hline
-Xmx512m & Límite de memoria del \textit{heap} (ajuste según RAM disponible). \\
\hline
\rowcolor{grisTecnico}
-XX:+UseContainerSupport & Respeta los límites de memoria del contenedor. \\
\hline
\seqsplit{-Djava.net.preferIPv4Stack=true} & Fuerza IPv4 (ya fijado por el \textit{plugin} de Spring Boot). \\
\hline
\rowcolor{grisTecnico}
-Dspring.profiles.active=prod & Activa el perfil de producción. \\
\hline
\end{tabularx}\normalfont\normalsize%
\label{tab:jvm}

\par\vspace{0.2cm}
\begin{lstlisting}[language=bash]
# Ejecucion con parametros de JVM ajustados
$ java -Xmx512m -XX:+UseContainerSupport \
    -jar target/ethoshub-0.0.1-SNAPSHOT.jar --spring.profiles.active=prod
\end{lstlisting}

En cuanto a observabilidad, Spring Boot Actuator expone, además de \comando{/health},
\textit{endpoints} de información y métricas que pueden integrarse con un sistema de
monitorización. La tabla~\ref{tab:obs} los resume.\par

\vspace{0.2cm}
\noindent
\renewcommand{\arraystretch}{1.3}
\fontsize{8.5}{10.2}\selectfont\sffamily
\begin{tabularx}{\dimexpr\textwidth-2pt\relax}{|>{\ttfamily\color{colorPrimario}}l|X|}
\hline
\rowcolor{headerTabla}
\sffamily\textbf{Endpoint} & \sffamily\textbf{Información} \\
\hline
/actuator/health & Estado de salud (UP/DOWN) del servicio y sus dependencias. \\
\hline
\rowcolor{grisTecnico}
/actuator/info & Metadatos de la aplicación. \\
\hline
/v3/api-docs & Especificación OpenAPI de la API. \\
\hline
\end{tabularx}\normalfont\normalsize%
\label{tab:obs}

% -----------------------------------------------------------------------------
\section{Anexo K — Matriz de Compatibilidad de Versiones}%
\label{sec:anexo-compat}

La tabla~\ref{tab:compat} fija las versiones probadas de cada componente. Desviarse de ellas
puede funcionar, pero queda fuera de la configuración validada.\par

\vspace{0.2cm}
\noindent
\renewcommand{\arraystretch}{1.3}
\fontsize{8.5}{10.2}\selectfont\sffamily
\begin{tabularx}{\dimexpr\textwidth-2pt\relax}{|>{\bfseries\color{colorPrimario}}l|l|X|}
\hline
\rowcolor{headerTabla}
\textbf{Componente} & \textbf{Versión validada} & \textbf{Notas} \\
\hline
Java (JDK) & 17 & Mínimo obligatorio; el \textit{build} aborta con versiones menores. \\
\hline
\rowcolor{grisTecnico}
Spring Boot & 4.0.6 & Heredado por \texttt{spring-boot-starter-parent}. \\
\hline
Node.js & 20 LTS & Compatible con la línea 25. \\
\hline
\rowcolor{grisTecnico}
React & 18.3 & SPA principal. \\
\hline
TypeScript & 5.3 & Modo \texttt{strict}. \\
\hline
\rowcolor{grisTecnico}
Vite & 5.1 & \textit{Bundler} y servidor de desarrollo. \\
\hline
PostgreSQL & 15 & Gestionado por Supabase. \\
\hline
\rowcolor{grisTecnico}
springdoc-openapi & 2.7.0 & Swagger UI / OpenAPI. \\
\hline
Tectonic & 0.15.0 & Motor LaTeX del CV Studio (en la imagen Docker). \\
\hline
\rowcolor{grisTecnico}
Docker Engine & 24+ & Con Compose v2. \\
\hline
\end{tabularx}\normalfont\normalsize%
\label{tab:compat}

% -----------------------------------------------------------------------------
\section{Anexo L — Mapa de Logs y Códigos de Estado}%
\label{sec:anexo-logs}

La tabla~\ref{tab:codigos} relaciona los códigos HTTP que el operador encontrará durante la
verificación con su interpretación y la acción sugerida.\par

\vspace{0.2cm}
\noindent
\renewcommand{\arraystretch}{1.3}
\fontsize{8.5}{10.2}\selectfont\sffamily
\begin{tabularx}{\dimexpr\textwidth-2pt\relax}{|c|>{\raggedright\arraybackslash}p{4.4cm}|X|}
\hline
\rowcolor{headerTabla}
\textbf{Código} & \textbf{Significado} & \textbf{Acción sugerida} \\
\hline
200 & Operación correcta. & Ninguna. \\
\hline
\rowcolor{grisTecnico}
401 & No autenticado (sin JWT o inválido). & Revisar token / JWKS (\ref{subsec:jwt}). \\
\hline
403 & Autenticado sin permisos. & Revisar rol del usuario. \\
\hline
\rowcolor{grisTecnico}
429 & Cuota de IA agotada (Gemini). & Reintentar más tarde (\ref{sec:degradacion-ia}). \\
\hline
502 & Error de pasarela hacia la IA. & Revisar conectividad a Gemini. \\
\hline
\rowcolor{grisTecnico}
503 & Servicio no disponible (BD o IA caída). & Revisar Supabase / egreso de red. \\
\hline
\end{tabularx}\normalfont\normalsize%
\label{tab:codigos}

\par\vspace{0.3cm}
Para soporte durante el despliegue, contacte al equipo de Byte Busters S.R.L. en
\ruta{contacto@bytebusters.com}. Este manual, identificado como
\ruta{MI-INST-ETHOSHUB-2.0.0}, constituye la referencia autoritativa del procedimiento
de instalación de la plataforma EthosHub.\par
